\documentclass[aspectratio=169,10pt]{beamer}

% Shared Beamer setup for Fall 2026 course slides.
\mode<presentation>{
  \usetheme{Madrid}
  \usecolortheme{default}
}

\definecolor{CalPolyGreen}{HTML}{154734}
\definecolor{CalPolyGold}{HTML}{BD8B13}
\definecolor{DeckInk}{HTML}{1F2933}
\definecolor{DeckMuted}{HTML}{5F6B7A}

\setbeamercolor{structure}{fg=CalPolyGreen}
\setbeamercolor{palette primary}{bg=CalPolyGreen,fg=white}
\setbeamercolor{palette secondary}{bg=DeckInk,fg=white}
\setbeamercolor{palette tertiary}{bg=CalPolyGold,fg=DeckInk}
\setbeamercolor{frametitle}{bg=CalPolyGreen,fg=white}
\setbeamercolor{title}{fg=CalPolyGreen}
\setbeamercolor{normal text}{fg=DeckInk,bg=white}
\setbeamercolor{block title}{bg=CalPolyGreen,fg=white}
\setbeamercolor{block body}{bg=black!3,fg=DeckInk}

\setbeamertemplate{navigation symbols}{}
\setbeamertemplate{footline}[frame number]
\setbeamertemplate{itemize item}{\color{CalPolyGold}$\blacktriangleright$}
\setbeamertemplate{itemize subitem}{\color{CalPolyGreen}$\bullet$}

\newcommand{\CourseNumber}{Course}
\newcommand{\CourseTitle}{Course Title}
\newcommand{\LectureNumber}{0}
\newcommand{\LectureTitle}{Lecture Title}
\newcommand{\LectureDate}{Date}
\newcommand{\CourseUnit}{Unit}

\newcommand{\coursenumber}[1]{\renewcommand{\CourseNumber}{#1}}
\newcommand{\coursetitle}[1]{\renewcommand{\CourseTitle}{#1}}
\newcommand{\lecturenumber}[1]{\renewcommand{\LectureNumber}{#1}}
\newcommand{\lecturetitle}[1]{\renewcommand{\LectureTitle}{#1}}
\newcommand{\lecturedate}[1]{\renewcommand{\LectureDate}{#1}}
\newcommand{\courseunit}[1]{\renewcommand{\CourseUnit}{#1}}

\author{Kirk Duran}
\institute{California Polytechnic State University, San Luis Obispo}
\date{\LectureDate}

\newcommand{\makedecktitle}{
  \begin{frame}[plain]
    \vspace{0.25in}
    {\usebeamerfont{title}\color{CalPolyGreen}\LectureTitle\par}
    \vspace{0.18in}
    {\large \CourseNumber{}: \CourseTitle\par}
    \vspace{0.08in}
    {\large Lecture \LectureNumber{} -- \LectureDate\par}
    \vfill
    {\small Kirk Duran \textbar{} Cal Poly, San Luis Obispo\par}
  \end{frame}
}

\newcommand{\placeholdernote}{
  \begin{block}{Placeholder}
    Replace these bullets with the final lecture material, examples, and in-class activities.
  \end{block}
}

\AtBeginSection[]{
  \begin{frame}{Roadmap}
    \tableofcontents[currentsection]
  \end{frame}
}


\graphicspath{{../assets/lecture-03/}}

% If an image file is not yet available, the slide displays a labeled
% placeholder using the intended filename. Replace the placeholder by placing
% the matching file in ../assets/lecture-03/.
\newcommand{\lectureimage}[2][1.75in]{%
  \IfFileExists{../assets/lecture-03/#2}{%
    \includegraphics[width=\linewidth,height=#1,keepaspectratio]{#2}%
  }{%
    \fbox{%
      \parbox[c][#1][c]{0.94\linewidth}{%
        \centering
        \small\textbf{Image Placeholder}\\[0.08in]
        \scriptsize\texttt{\detokenize{#2}}
      }%
    }%
  }%
}

\setbeamersize{text margin left=0.42in,text margin right=0.42in}

\coursenumber{CS 1001}
\coursetitle{Introduction to Computing}
\lecturenumber{3}
\lecturetitle{Python Expressions and Evaluation}
\lecturedate{August 28, 2026}
\courseunit{1. Computing and Program Execution}

\begin{document}

\makedecktitle

\begin{frame}{Today}

  \begin{itemize}

    \item Why programming languages let us work above machine-level details.

    \item How low-level and high-level languages differ.

    \item How compilation and interpretation connect source code to execution.

    \item Why Python is widely used and appropriate for this course.

    \item How Python evaluates one self-contained expression to produce one value.

  \end{itemize}

  \begin{keyidea}

    Today we evaluate expressions without introducing variables or persistent program state.

  \end{keyidea}

\end{frame}

\sectionframe{Part 1: Moving Away from Bare Metal}

\begin{frame}{Programming Languages Bridge Human Ideas and Machine Execution}

  \begin{columns}[T,onlytextwidth]

    \begin{column}{0.53\textwidth}

      \begin{cpdefinition}

        A programming language is a formal notation for expressing computations that a computer can execute.

      \end{cpdefinition}

      \begin{itemize}

        \item Its \textbf{syntax} specifies which forms are valid.

        \item Its \textbf{semantics} specifies what valid forms mean.

        \item Software translates or processes the notation so hardware can perform the requested computation.

      \end{itemize}

    \end{column}

    \begin{column}{0.41\textwidth}

      % Search direction: human-readable code bridging to CPU instructions.
      \lectureimage[1.95in]{4.png}

    \end{column}

  \end{columns}

\end{frame}

\begin{frame}{Low-Level Languages Stay Close to Hardware}

  \begin{columns}[T,onlytextwidth]

    \begin{column}{0.52\textwidth}

      \begin{block}{Low-level languages}

        \begin{itemize}

          \item Closely reflect processor operations.

          \item Expose memory and hardware details.

          \item Offer precise control but require more implementation detail.

        \end{itemize}

      \end{block}

      \begin{keyidea}

        Low-level code gives the programmer more direct control over what the machine does, but that control comes with more machine-specific detail.

      \end{keyidea}

    \end{column}

    \begin{column}{0.42\textwidth}

      % Search direction: machine code and assembly instructions near CPU hardware.
      \lectureimage[2.20in]{5.jpg}

    \end{column}

  \end{columns}

\end{frame}

\begin{frame}{High-Level Languages Hide Machine Detail}

  \begin{columns}[T,onlytextwidth]

    \begin{column}{0.52\textwidth}

      \begin{block}{High-level languages}

        \begin{itemize}

          \item Express computations using human-oriented abstractions.

          \item Hide many processor-specific details.

          \item Improve readability, portability, and development speed.

        \end{itemize}

      \end{block}

      \begin{keyidea}

        High-level languages let us write closer to the problem we are solving and farther from the processor instructions that eventually run.

      \end{keyidea}

    \end{column}

    \begin{column}{0.42\textwidth}

      % Search direction: abstraction ladder from machine code through assembly to Python.
      \lectureimage[2.20in]{6.jpg}

    \end{column}

  \end{columns}

\end{frame}

\begin{frame}{Compiled Languages Translate Before Execution}

  \begin{columns}[T,onlytextwidth]

    \begin{column}{0.52\textwidth}

      \begin{block}{Compilation}

        A compiler translates source code into another executable form before that form runs.

      \end{block}

      \begin{itemize}

        \item The compiler checks and translates source code ahead of execution.

        \item The output may be machine code, bytecode, or another lower-level form.

        \item The translated form can often run many times without recompiling from scratch.

      \end{itemize}

    \end{column}

    \begin{column}{0.42\textwidth}

      % Search direction: compiler turning source code into executable instructions.
      \lectureimage[2.15in]{7.jpg}

    \end{column}

  \end{columns}

\end{frame}

\begin{frame}[shrink=10]{Interpreted Languages Process Code During Execution}

  \begin{columns}[T,onlytextwidth]

    \begin{column}{0.52\textwidth}

      \begin{block}{Interpretation}

        An interpreter or runtime processes a program and performs its operations during execution.

      \end{block}

      \begin{itemize}

        \item The runtime reads program forms and carries out their meaning as the program runs.

        \item This can support fast experimentation and interactive use.

        \item Many real systems blend strategies; CPython compiles source to bytecode and then executes that bytecode.

      \end{itemize}

    \end{column}

    \begin{column}{0.42\textwidth}

      % Search direction: compiler and interpreter paths converging on machine execution.
      \lectureimage[1.55in]{8.png}

    \end{column}

  \end{columns}

  \begin{keyidea}

    Compiled and interpreted describe implementation strategies, not perfectly separate categories of languages.

  \end{keyidea}

\end{frame}

\begin{frame}{Python Keeps Our Attention on the Computation}

  \begin{columns}[T,onlytextwidth]

    \begin{column}{0.53\textwidth}

      \begin{itemize}

        \item Python is a high-level, general-purpose programming language.

        \item Its syntax emphasizes readability and supports rapid experimentation.

        \item It runs across major operating systems and has a large library ecosystem.

        \item It is used in automation, scientific computing, data analysis, artificial intelligence, web systems, and education.

      \end{itemize}

      \begin{exampleblockcp}

        In CPython, source code is compiled to bytecode, which the Python virtual machine executes.

      \end{exampleblockcp}

    \end{column}

    \begin{column}{0.41\textwidth}

      % Search direction: Python applications in science, automation, data, AI, and web development.
      \lectureimage[2.18in]{9.jpg}

    \end{column}

  \end{columns}

\end{frame}

\sectionframe{Part 2: A Brief History of Programming Languages}

\begin{frame}{Early Programming Stayed Close to the Machine}

  \begin{columns}[T,onlytextwidth]

    \begin{column}{0.54\textwidth}

      \begin{itemize}

        \item Early electronic computers were programmed through switches, wiring, or numerical instruction codes.

        \item Machine language encoded processor instructions directly as numerical patterns.

        \item Assembly language replaced many numerical patterns with symbolic instruction names.

        \item Programs remained closely tied to a particular processor architecture.

      \end{itemize}

    \end{column}

    \begin{column}{0.40\textwidth}

      % Search direction: early computer programming, punched cards, and assembly notation.
      \lectureimage[2.15in]{11.jpg}

    \end{column}

  \end{columns}

\end{frame}

\begin{frame}{Early High-Level Languages Described Problem Domains}

  \begin{columns}[T,onlytextwidth]

    \begin{column}{0.54\textwidth}

      \begin{itemize}

        \item \textbf{FORTRAN} supported scientific and numerical computing.

        \item \textbf{COBOL} supported business data processing.

        \item \textbf{Lisp} supported symbolic computation and early artificial intelligence research.

        \item Compilers translated these higher-level descriptions into lower-level instructions.

      \end{itemize}

      \begin{keyidea}

        Programmers increasingly described the problem to solve rather than every processor operation.

      \end{keyidea}

    \end{column}

    \begin{column}{0.40\textwidth}

      % Search direction: FORTRAN, COBOL, and Lisp historical montage.
      \lectureimage[2.15in]{12.jpeg}

    \end{column}

  \end{columns}

\end{frame}

\begin{frame}{Languages Added Structure, Portability, and New Models}

  \begin{columns}[T,onlytextwidth]

    \begin{column}{0.54\textwidth}

      \begin{itemize}

        \item \textbf{BASIC} made interactive programming more accessible.

        \item \textbf{C} combined portable systems programming with low-level control.

        \item \textbf{Pascal} emphasized structured programming and clear program organization.

        \item \textbf{Smalltalk} advanced object-oriented programming and graphical environments.

      \end{itemize}

    \end{column}

    \begin{column}{0.40\textwidth}

      % Search direction: BASIC, C, Pascal, and Smalltalk language timeline.
      \lectureimage[2.15in]{13.png}

    \end{column}

  \end{columns}

\end{frame}

\begin{frame}{Modern Languages Grow Through Their Ecosystems}

  \begin{columns}[T,onlytextwidth]

    \begin{column}{0.54\textwidth}

      \begin{itemize}

        \item Languages such as Python, Java, and JavaScript expanded with the internet and personal computing.

        \item Libraries package reusable solutions for common tasks.

        \item Development tools support testing, debugging, collaboration, and deployment.

        \item Communities, documentation, and application domains strongly influence adoption.

      \end{itemize}

    \end{column}

    \begin{column}{0.40\textwidth}

      % Search direction: modern programming languages surrounded by tools, libraries, and communities.
      \lectureimage[2.15in]{14.png}

    \end{column}

  \end{columns}

\end{frame}

\sectionframe{Part 3: Python Expressions Produce Values}

\begin{frame}[shrink=10]{A Literal Writes a Value Directly in Code}

  \begin{columns}[T,onlytextwidth]

    \begin{column}{0.48\textwidth}

      \begin{cpdefinition}

        A \textbf{literal} is a value written directly in source code.

      \end{cpdefinition}

      \begin{itemize}

        \item \texttt{42}

        \item \texttt{3.14}

        \item \texttt{"hello"}

        \item \texttt{True}

      \end{itemize}

    \end{column}

    \begin{column}{0.46\textwidth}

      % Search direction: simple values entering a Python evaluation pipeline.
      \lectureimage[1.55in]{16.jpeg}

    \end{column}

  \end{columns}

  \begin{keyidea}

    The written form represents a value; no name or storage location is needed for today's examples.

  \end{keyidea}

\end{frame}

\begin{frame}{An Expression Evaluates to Produce One Value}

  \begin{columns}[T,onlytextwidth]

    \begin{column}{0.50\textwidth}

      \begin{cpdefinition}

        An \textbf{expression} is code that Python evaluates to produce one value.

      \end{cpdefinition}

      \begin{itemize}

        \item \texttt{4 + 5} $\rightarrow$ \texttt{9}

        \item \texttt{3 * (8 - 2)} $\rightarrow$ \texttt{18}

        \item \texttt{"CS" + "1001"} $\rightarrow$ \texttt{"CS1001"}

      \end{itemize}

    \end{column}

    \begin{column}{0.44\textwidth}

      \begin{block}{Mental model}

        \centering
        values\\[0.06in]
        $\downarrow$\\[0.06in]
        operator\\[0.06in]
        $\downarrow$\\[0.06in]
        new value

      \end{block}

    \end{column}

  \end{columns}

\end{frame}

\begin{frame}{Operators Combine or Transform Values}

  \begin{columns}[T,onlytextwidth]

    \begin{column}{0.50\textwidth}

      \begin{itemize}

        \item \texttt{+} addition

        \item \texttt{-} subtraction

        \item \texttt{*} multiplication

        \item \texttt{/} division

      \end{itemize}

    \end{column}

    \begin{column}{0.44\textwidth}

      \begin{itemize}

        \item \texttt{//} floor division

        \item \texttt{\%} modulus

        \item \texttt{**} exponentiation

        \item \texttt{( )} explicit grouping

      \end{itemize}

    \end{column}

  \end{columns}

  \begin{exampleblockcp}

    Each operator receives one or more values and produces a new value for the surrounding expression.

  \end{exampleblockcp}

\end{frame}

\begin{frame}{Python Evaluates Nested Expressions from the Inside Out}

  \begin{center}

    {\Large \texttt{3 * (8 - 2)}}

  \end{center}

  \vspace{0.08in}

  \begin{enumerate}

    \item The parentheses identify the inner expression: \texttt{8 - 2}.

    \item That subexpression produces \texttt{6}.

    \item The surrounding expression becomes \texttt{3 * 6}.

    \item Multiplication produces the final value \texttt{18}.

  \end{enumerate}

  \begin{keyidea}

    A completed subexpression becomes a value used by the surrounding operation.

  \end{keyidea}

\end{frame}

\begin{frame}{Division Produces a Quotient}

  \begin{columns}[T,onlytextwidth]

    \begin{column}{0.50\textwidth}

      \begin{itemize}

        \item \texttt{7 + 2} $\rightarrow$ \texttt{9}

        \item \texttt{7 - 2} $\rightarrow$ \texttt{5}

        \item \texttt{7 * 2} $\rightarrow$ \texttt{14}

        \item \texttt{7 / 2} $\rightarrow$ \texttt{3.5}

      \end{itemize}

    \end{column}

    \begin{column}{0.44\textwidth}

      \begin{exampleblockcp}

        \texttt{8 / 4} produces \texttt{2.0}, not \texttt{2}.

      \end{exampleblockcp}

      \begin{keyidea}

        In Python, the \texttt{/} operator produces a floating-point quotient, even when the division is exact.

      \end{keyidea}

    \end{column}

  \end{columns}

\end{frame}

\begin{frame}{Floor Division Counts Complete Groups}

  \begin{columns}[T,onlytextwidth]

    \begin{column}{0.47\textwidth}

      \begin{center}

        {\LARGE \texttt{17 // 5} $\rightarrow$ \texttt{3}}

      \end{center}

      \begin{itemize}

        \item Seventeen items contain three complete groups of five.

        \item The two remaining items do not form another complete group.

        \item For now, we use nonnegative values; negative operands require a more precise discussion of ``floor.''

      \end{itemize}

    \end{column}

    \begin{column}{0.47\textwidth}

      % Search direction: 17 objects arranged as three complete groups of five plus two extras.
      \lectureimage[2.20in]{21.png}

    \end{column}

  \end{columns}

\end{frame}

\begin{frame}{Modulus Produces the Remainder}

  \begin{columns}[T,onlytextwidth]

    \begin{column}{0.47\textwidth}

      \begin{center}

        {\LARGE \texttt{17 \% 5} $\rightarrow$ \texttt{2}}

      \end{center}

      \begin{itemize}

        \item Three complete groups of five use fifteen items.

        \item Two items remain.

        \item Modulus is useful for cycles, divisibility, and extracting final digits.

      \end{itemize}

      \begin{exampleblockcp}

        \texttt{123 \% 10} $\rightarrow$ \texttt{3}

      \end{exampleblockcp}

    \end{column}

    \begin{column}{0.47\textwidth}

      % Search direction: 17 objects in groups of five with two remainders highlighted.
      \lectureimage[2.20in]{22.jpg}

    \end{column}

  \end{columns}

\end{frame}

\begin{frame}{Exponentiation Repeats Multiplication}

  \begin{columns}[T,onlytextwidth]

    \begin{column}{0.49\textwidth}

      \begin{center}

        {\LARGE \texttt{2 ** 3} $\rightarrow$ \texttt{8}}

      \end{center}

      \begin{block}{Interpretation}

        \texttt{2 ** 3} means three factors of two:

        \begin{center}
          $2 \times 2 \times 2 = 8$
        \end{center}

      \end{block}

    \end{column}

    \begin{column}{0.45\textwidth}

      \begin{itemize}

        \item Python uses \texttt{**} for exponentiation.

        \item The symbol \texttt{\string^} does not mean exponentiation in Python.

        \item Exponentiation has higher precedence than multiplication, division, addition, and subtraction.

      \end{itemize}

    \end{column}

  \end{columns}

\end{frame}

\sectionframe{Part 4: Order of Operations}

\begin{frame}{Precedence Determines Which Operation Acts First}

  \begin{center}

    \begin{tabular}{c l l}
      \textbf{Priority} & \textbf{Operators} & \textbf{Rule} \\
      \hline
      1 & \texttt{( )} & Innermost grouping first \\
      2 & \texttt{**} & Right to left \\
      3 & \texttt{*}, \texttt{/}, \texttt{//}, \texttt{\%} & Left to right \\
      4 & \texttt{+}, \texttt{-} & Left to right \\
    \end{tabular}

  \end{center}

  \vspace{0.12in}

  \begin{keyidea}

    PEMDAS is a mnemonic. The table states the Python rules more precisely: multiplication and division share a level, as do addition and subtraction.

  \end{keyidea}

\end{frame}

\begin{frame}{Associativity Resolves Operators at the Same Level}

  \begin{columns}[T,onlytextwidth]

    \begin{column}{0.47\textwidth}

      \begin{block}{Usually left to right}

        \texttt{20 / 5 * 2}

        \vspace{0.08in}

        \texttt{4.0 * 2} $\rightarrow$ \texttt{8.0}

      \end{block}

      Multiplication does not automatically occur before division; the two operators have equal precedence.

    \end{column}

    \begin{column}{0.47\textwidth}

      \begin{block}{Exponentiation: right to left}

        \texttt{2 ** 3 ** 2}

        \vspace{0.08in}

        \texttt{2 ** (3 ** 2)} $\rightarrow$ \texttt{512}

      \end{block}

      Exponentiation is the important exception in today's operator set.

    \end{column}

  \end{columns}

\end{frame}

\begin{frame}{Parentheses Override the Default Grouping}

  \begin{columns}[T,onlytextwidth]

    \begin{column}{0.47\textwidth}

      \begin{block}{Default precedence}

        \begin{center}
          {\Large \texttt{2 + 3 * 4}}\\[0.10in]
          \texttt{2 + 12} $\rightarrow$ \texttt{14}
        \end{center}

      \end{block}

    \end{column}

    \begin{column}{0.47\textwidth}

      \begin{block}{Explicit grouping}

        \begin{center}
          {\Large \texttt{(2 + 3) * 4}}\\[0.10in]
          \texttt{5 * 4} $\rightarrow$ \texttt{20}
        \end{center}

      \end{block}

    \end{column}

  \end{columns}

  \begin{keyidea}

    Parentheses can change a result and can make the intended grouping easier to read.

  \end{keyidea}

\end{frame}

\begin{frame}[shrink=10]{Reminder: The Processor Executes One Step at a Time}

  \begin{columns}[T,onlytextwidth]

    \begin{column}{0.52\textwidth}

      \begin{block}{Von Neumann architecture}

        A stored-program computer keeps instructions and data in memory. The CPU repeatedly fetches instructions, decodes them, executes them, and stores results.

      \end{block}

      \begin{block}{Processing cycle}

        Fetch an instruction from memory, decode the operation, execute it using CPU hardware, then move to the next instruction.

      \end{block}

    \end{column}

    \begin{column}{0.42\textwidth}

      \begin{keyidea}

        Arithmetic operators eventually become processor operations carried out by the arithmetic logic unit (ALU).

      \end{keyidea}

      \begin{exampleblockcp}

        For today, ignore variables and stored program state. We are only tracing how one expression produces one value.

      \end{exampleblockcp}
    \end{column}

  \end{columns}

\end{frame}

\sectionframe{Part 5: Predict, Verify, and Explain}

\begin{frame}{Predict Each Value Before Running Python}

  \begin{columns}[T,onlytextwidth]

    \begin{column}{0.44\textwidth}

      \begin{enumerate}

        \item \texttt{2 + 3 * 4}

        \item \texttt{(2 + 3) * 4}

        \item \texttt{17 // 5}

        \item \texttt{17 \% 5}

        \item \texttt{2 ** 3 ** 2}

      \end{enumerate}

    \end{column}

    \begin{column}{0.50\textwidth}

      \begin{activity}

        For each expression, record:

        \begin{enumerate}

          \item the final value,

          \item the first operation performed, and

          \item whether parentheses alter the result.

        \end{enumerate}

      \end{activity}

    \end{column}

  \end{columns}

\end{frame}

\begin{frame}[shrink=10]{Use Colab to Test a Reasoned Prediction}

  \begin{columns}[T,onlytextwidth]

    \begin{column}{0.47\textwidth}

      \begin{enumerate}

        \item Predict the final value.

        \item Identify the expected evaluation order.

        \item Enter one expression in one Colab cell.

        \item Run the cell and compare the result.

        \item Explain any difference between the prediction and Python's result.

      \end{enumerate}

    \end{column}

    \begin{column}{0.47\textwidth}

      % Search direction: Google Colab cell evaluating 2 + 3 * 4 and displaying 14.
      \lectureimage[1.55in]{31.jpeg}

    \end{column}

  \end{columns}

  \begin{keyidea}

    Colab verifies a prediction; it does not replace an explanation.

  \end{keyidea}

\end{frame}

\begin{frame}{Evaluation and Display Are Different Actions}

  \begin{columns}[T,onlytextwidth]

    \begin{column}{0.47\textwidth}

      \begin{block}{Evaluate an expression}

        \begin{center}
          \texttt{2 + 3}
        \end{center}

        The expression evaluates to \texttt{5}. Colab automatically displays the value of the final expression in a cell.

      \end{block}

    \end{column}

    \begin{column}{0.47\textwidth}

      \begin{block}{Request display explicitly}

        \begin{center}
          \texttt{print(2 + 3)}
        \end{center}

        Python first evaluates \texttt{2 + 3}; then \texttt{print()} displays the resulting value.

      \end{block}

    \end{column}

  \end{columns}

  \begin{exampleblockcp}

    In a script, a bare expression usually produces no visible output. Evaluation still occurs.

  \end{exampleblockcp}

\end{frame}

\begin{frame}{One Expression Produces One Explainable Value}

  \begin{activity}

    Evaluate the following expression without running Python first:

    \begin{center}

      {\LARGE \texttt{17 \% 5 + 2 ** 3 * 2}}

    \end{center}

    Report the final value and annotate the order of evaluation.

  \end{activity}

  \begin{keyidea}

    Evaluate exponentiation first, then multiplication and division-level operators, then addition and subtraction. Parentheses override the default order.

  \end{keyidea}

\end{frame}

\begin{frame}{Next Time}

  \begin{center}

    {\Huge Naming and Retaining Values}

    \vspace{0.20in}

    {\Large Variables and Assignment}

  \end{center}

\end{frame}

\end{document}
